\begin{abstract}
\noindent
Retrieval-Augmented Generation (RAG) systems can provide information from
external knowledge sources, but they still face difficulties with legal
questions that require connecting multiple rules, identifying reasoning
paths, and examining counterfactual assumptions. This study presents a
CausalRAG pipeline for question answering in Vietnamese criminal law, in
which legal rules are normalized into condition--effect relations and
organized as a heterogeneous legal causal graph. The system combines a
BGE-M3-based vector memory, multi-hop causal-path retrieval, query-aware
claim verification, and the generation of answers grounded in statutory
provisions.

The central verification component of the system is a Legal Structural
Causal Model (LegalSCM), a deterministic normative reasoning model with a
three-valued state space. LegalSCM validates factual causal paths and
performs a hard intervention by assigning an intermediate event the
\(\mathrm{FALSE}\) state, severing its corresponding incoming mechanisms,
and re-inferring the outcome. A baseline based on node deletion and graph
reachability is used to conduct a paired ablation.

Experiments are conducted on 2,884 rules, 1,814 events, and a synthetic
evaluation set comprising 250 multi-hop questions. The system achieves a
Rule Recall@5 of 69.93\%, an Event Recall@5 of 74.66\%, a verification
accuracy of 92.00\%, an Answer Token F1 of 68.69\%, and a Citation F1 of
62.86\%. Top-1 Exact Path reaches 35.20\%, whereas Oracle Exact Path reaches
71.20\%, indicating that path ranking remains the principal bottleneck. The
paired ablation reveals no difference in quality between LegalSCM and path
ablation on the current benchmark; however, LegalSCM provides clearer
intervention semantics and more detailed reasoning traces, at an average
processing time approximately 3.23 times higher. The results demonstrate
the potential of combining graph-based retrieval, structured verification,
and grounded answer generation, while also highlighting the need for
counterfactual benchmarks that are more sensitive to intervention
mechanisms.
\end{abstract}

\vspace{0.35em}
\noindent\textbf{Keywords:}
Retrieval-Augmented Generation; CausalRAG; causal graph;
structural causal model; counterfactual reasoning;
legal question answering; Vietnamese Penal Code.
\vspace{0.6em}
